% ============================================================
%  Chapter 3: Related Work
% ============================================================
%
%  TARGET: ~600 words (DOWN FROM ~2,600 in current draft).
%  ~4x compression. Most of the literature survey has to go.
%
%  STRATEGY: this is not a literature review. It's a positioning
%  argument. The reader needs to know (a) what families of
%  learned reconstruction exist, (b) where TFPnP sits among them,
%  (c) what TFPnP claims and where its gaps are.
%
% ============================================================
%
%  KEEP from current draft:
%   - The PnP parameter-tuning paragraph (compressed into Ch.2
%     §2.3 already; just one cross-reference here)
%   - TFPnP MDP formulation in brief
%   - TFPnP key innovations as bullet sentences
%   - "Assumptions and Limitations" paragraph -- this is where
%     you show independent thinking and it's high-value per word
%
%  CUT from current draft (saves ~2000 words):
%   - Section 3.1 (Classical iterative) -- covered in Ch.2, drop entirely
%   - Section 3.2 (Variational hierarchy) -- compress to one sentence
%   - Section 3.3 (Proximal algorithms) -- the equations live in Ch.2; drop
%   - Section 3.4 (PnP variants list HQS/PGM/APGM/RED/D-AMP) -- one sentence
%   - Section 3.5 (LPD full description) -- one sentence as baseline only
%   - Section 3.6 (DL for CT taxonomy) -- one sentence per family
%   - Paper authors + JMLR vol detail -- not needed
%   - Full reward function and discounted-return equation -- one line
%   - Comparison subsections (vs PnP, vs supervised, vs unrolling) --
%     collapse into one "positioning" paragraph
%   - "What the Paper Does Not Claim" -- merge with limitations
%
% ============================================================
%
%  TARGET STRUCTURE (3 sections, no subsections):
%
%  3.1 Learned reconstruction landscape       ~150 w
%       ¶ Two-axis taxonomy: (i) whether the forward operator A
%         enters at inference, (ii) supervised pixel-wise vs
%         RL-/self-supervised. Name the families with one citation
%         each: post-processing (FBPconv, RED-CNN), unrolled
%         (LPD, RPGD), PnP (Venkatakrishnan, IRCNN). Mention LPD
%         is the SOTA supervised benchmark for sparse-view CT
%         (table later).
%
%  3.2 TFPnP                                  ~350 w
%       ¶ One paragraph: TFPnP (Wei et al., JMLR 2022) frames the
%         PnP schedule problem as a Markov decision process. State =
%         current (x,z,u) plus iteration count and noise level;
%         action = (sigma_k, mu_k, terminate?); reward =
%         per-step PSNR improvement. Cite RL background lightly.
%       ¶ Second paragraph: training combines (i) REINFORCE for
%         the discrete termination head pi_1, (ii) deterministic
%         policy gradient through the differentiable PnP iterations
%         for the continuous head pi_2, and (iii) a critic V.
%         This is the "mixed model-free/model-based" framing.
%       ¶ Third paragraph: paper's headline CT result. Insert
%         Table 3.1 here (the paper's own numbers from Table 5);
%         this same table is referenced again in Ch.5 for
%         comparison with this project's results.
%
%  3.3 Assumptions, limitations, and position of this work   ~100 w
%       ¶ Three crisp limitations of the paper:
%           (i) cross-domain denoiser (trained on PASCAL VOC,
%               applied to CT) -- this work tests the in-domain
%               variant.
%           (ii) TorchRadon backend, only 2D parallel-beam --
%               this work uses LION/tomosipo.
%           (iii) IRCNN denoiser -- this work uses DRUNet,
%               testing the off-the-shelf generalisation claim.
%       ¶ One sentence positioning: these three differences
%         define the contributions enumerated in Chapter 1.
%
% ============================================================
%  FIGURES / TABLES
%
%  Table 3.1: Paper Table 5 reproduced (PSNR for FBP / TV /
%             FBPconv / RED-CNN / LPD / RPGD / TFPnP at 5/7.5/10%).
%             KEEP from current draft. Caption ~20 words.
%             This table is forward-referenced from Chapter 5.
%
%  NO FIGURES in this chapter. The PnP-ADMM schematic figure
%  hinted at in the original draft (Fig 3.1) is CUT for space.
% ============================================================

\chapter{Related Work}
\label{chap:related-work}

% -----------------------------------------------------------
\section{Learned Reconstruction: A Brief Landscape}
\label{sec:learned-landscape}
% -----------------------------------------------------------

% TODO [single paragraph, ~150 w]
% Taxonomise learned CT reconstruction along two axes:
%   (a) Image-domain post-processing (FBPconv \citep{Jin2017},
%       RED-CNN \citep{Chen2017RED}). Fast at inference but no
%       data-consistency enforcement.
%   (b) Unrolled algorithms (LPD \citep{AdlerOktem2018}, RPGD
%       \citep{Gupta2018}). State-of-the-art supervised
%       performance; one trained model per geometry/noise level.
%   (c) Plug-and-Play (Venkatakrishnan2013 etc.; covered in Ch.2).
%       Modular, no paired training required, but sensitive to
%       hyperparameter schedules.
% End with one sentence stating that TFPnP addresses the schedule
% sensitivity problem in (c) via reinforcement learning.


% -----------------------------------------------------------
\section{TFPnP: Reinforcement-Learned Parameter Schedules}
\label{sec:tfpnp}
% -----------------------------------------------------------

% TODO ¶1 [MDP formulation, ~120 w]
% TFPnP (Wei et al., 2022) frames per-iteration PnP parameter
% selection as a Markov decision process:
%   - State: (x^k, z^k, u^k) + iteration index + noise level.
%   - Action: (sigma_k, mu_k) for the next m PnP-ADMM steps,
%       plus a discrete termination flag.
%   - Reward: r_t = PSNR(x^{t+1}) - PSNR(x^t) - eta (default
%       eta = 0.05 penalises additional iterations).
% A single policy is trained across multiple noise levels.

% TODO ¶2 [training algorithm, ~120 w]
% The policy has two heads: a discrete termination head pi_1
% trained by REINFORCE, and a continuous (sigma, mu) head pi_2
% trained by deterministic policy gradient through the
% differentiable PnP iterations. A critic V regresses to the
% empirical return. The combination is called "mixed model-free /
% model-based" in the paper. Mention that pi_2 training is
% sensitive (this is foreshadowing for Ch.4 hyperparameter
% choices and Ch.6 discussion of training instability).

% TODO ¶3 [CT results table, ~50 w]
% Reference Table 3.1: TFPnP achieves the best PSNR at all three
% noise levels, beating RPGD (the strongest unrolled baseline) by
% 0.21-0.34 dB. Single sentence; the table does the work.

\begin{table}[htbp]
    \centering
    \caption{Sparse-view CT PSNR (dB) from \citet{Wei2022TFPnP},
             Table 5. The paper's headline result against which
             this project's reproduction is compared in
             \cref{chap:results}.}
    \label{tab:paper-ct-results}
    \begin{tabular}{lccc}
        \toprule
        \textbf{Method} & \textbf{5\%} & \textbf{7.5\%} & \textbf{10\%} \\
        \midrule
        FBP \citep{Herman2009}          & 18.37 & 16.16 & 14.32 \\
        TV \citep{SidkyPan2008}         & 21.47 & 19.17 & 16.95 \\
        FBPconv \citep{Jin2017}         & 23.00 & 22.94 & 22.86 \\
        RED-CNN \citep{Chen2017RED}     & 23.45 & 23.11 & 22.71 \\
        LPD \citep{AdlerOktem2018}      & 23.86 & 23.44 & 23.00 \\
        RPGD \citep{Gupta2018}          & 24.29 & 23.82 & 23.38 \\
        \midrule
        \textbf{TFPnP}                  & \textbf{24.50} & \textbf{24.23} & \textbf{23.72} \\
        \bottomrule
    \end{tabular}
\end{table}


% -----------------------------------------------------------
\section{Position of This Work}
\label{sec:positioning}
% -----------------------------------------------------------

% TODO [~150 w, ends the chapter and bridges to Methods]
% Three deliberate differences from the paper, each motivating
% an experiment in Chapter 5:
%
%   1. IN-DOMAIN denoiser training. The paper trains its denoiser
%      on PASCAL VOC natural images and tests cross-domain on CT.
%      This work trains and tests on LIDC-IDRI -- testing the
%      upper bound of TFPnP under domain match. This complicates
%      direct numerical comparison with the paper (a caveat
%      repeated in Ch.6).
%
%   2. OFF-THE-SHELF DRUNet denoiser (vs IRCNN). DRUNet is a
%      stronger and more widely used denoiser. Using it here
%      tests the paper's central generalisation claim: "any
%      off-the-shelf denoiser works".
%
%   3. LION/tomosipo backend (vs TorchRadon). Engineering
%      contribution that makes the TFPnP framework available to
%      the wider Cambridge CIA group, with cone-beam ready for
%      future work.

% End with single sentence: "Chapter 4 describes how these three
% differences were realised in code; Chapter 5 reports their
% empirical effect."